\section{从第一家公司到阿里：移动互联网入口的前史}

本节先看祝铭明的第一段创业，因为它解释了 Rokid 后来为什么总是盯着“下一代入口”。2007 年 iPhone 发布后，他判断智能机和移动互联网会成为新平台，于是做了一家移动操作系统公司，希望把类似 iPhone 的软件体系迁移到不同硬件平台。后来阿里准备进入移动互联网，需要一个基础抓手，最终以约一千万美元收购了这家公司。

这个故事的关键不是“卖给阿里”，而是三点经验。第一，入口切换往往出现在硬件和软件体系同时变化的时刻；第二，创业公司可能在最危急的时候被大公司看见；第三，平台级机会要求创始人能忍受早期极高的不确定性。祝铭明回忆，公司最糟糕时账上只有不到四千块，两周后就要发工资，甚至准备让员工把电脑搬回去抵工资。

\begin{figure}[H]
\centering
\includegraphics[width=0.96\textwidth]{rokid-founder-timeline.png}
\caption{Rokid 创业时间线：从阿里收购、M Lab 到 AI/AR 硬件十一年。自制概念图，依据 00:02:36--00:31:41 对谈内容整理。}
\end{figure}

\begin{knowledgebox}{读图：第一段创业不是插曲，而是入口判断训练}
图中从 First Startup 到 Alibaba，再到 M Lab 和 Rokid，展示的是同一类判断：新计算平台出现时，底层系统、入口硬件和应用生态会重新洗牌。祝铭明从移动 OS 进入阿里，再出来做 AI/AR，本质上一直在找下一代入口。
\end{knowledgebox}

\subsection{危机时刻为什么重要}

本节把最早那段“账上不到四千块”的故事单独拿出来看。它不是励志段落，而是入口创业的现金流训练。祝铭明回忆，当时两周后就要发工资，甚至准备让员工把电脑搬回去抵工资；就在这个阶段，马云、蔡崇信和王坚博士把阿里收购的路径带了进来。创业公司被看见，并不总发生在最光鲜的时候，反而可能发生在方向被少数人理解、现金流已经接近极限的时候。

\begin{warningbox}{实践经验：方向正确不等于现金流安全}
一个创业项目即使押中了下一代入口，也可能在商业化、融资或大客户出现前耗尽现金。访谈里最早的移动 OS 公司和后来的 Rokid 都说明：前沿方向需要战略耐心，但公司每天仍要面对工资、供应商、库存和团队信心。
\end{warningbox}

\noindent\begin{tabularx}{\textwidth}{p{0.20\textwidth}p{0.34\textwidth}X}
\toprule
事件 & 表面故事 & 教学含义 \\
\midrule
2007 年 iPhone 发布 & 创始人看到智能机方向 & 新入口出现时，旧硬件和旧软件体系都会松动。 \\
做移动 OS & 把 iPhone 式软件体系迁移到多种机型 & 平台机会往往先表现为底层系统和兼容性问题。 \\
阿里收购 & 一千万美元股票交易 & 巨头要进入新平台时，需要基础能力和团队。 \\
现金流危机 & 账上钱不够发工资 & 入口判断再宏大，也必须穿过公司生存线。 \\
\bottomrule
\end{tabularx}

\subsection{阿里、吴妈与 M Lab}

本节解释阿里经历如何影响 Rokid。祝铭明进入阿里后，一开始参与移动互联网相关业务，后来吴泳铭提出要做 AI，成立 M Lab。那时深度学习刚起步，语音识别、拍照搜索、AIoT 和智能硬件都处在早期探索阶段。阿里想做 fundamental 的东西：平台、系统、基础能力，而不只是一个单点应用。

M Lab 的经验给 Rokid 留下两个长期影响。第一，AI 不是一个孤立功能，而是会寻找硬件入口。第二，大公司可以组织资源，但在某些前沿硬件方向上，创业公司更适合用高风险、高专注来先跑。吴泳铭后来也是 Rokid 的天使投资人，这说明这段关系不只是履历背景，也影响了后续资金和判断。

\begin{knowledgebox}{术语消化：入口、平台、产品}
\noindent\begin{tabularx}{\textwidth}{p{0.18\textwidth}p{0.34\textwidth}X}
\toprule
概念 & 含义 & 本期中的判断 \\
\midrule
产品 & 解决一个具体需求，收入来自单品销售或服务 & AI 音箱更接近产品，不一定是平台。 \\
入口 & 用户频繁进入数字世界的默认通道 & 手机是上一代入口，眼镜可能承担碎片化入口。 \\
平台 & 能承载开发者、应用和生态的基础层 & 智能眼镜若有足够使用时长和开发者，才可能平台化。 \\
\bottomrule
\end{tabularx}
\end{knowledgebox}

\begin{importantbox}{老师强调：fundamental 的东西会改变组织边界}
阿里当时吸收移动 OS 团队，是因为移动互联网需要底层抓手；吴泳铭推动 AI 和 M Lab，也是因为 AI 不能只被当作一个功能插件。Rokid 后来持续强调 OS、平台和生态，来自同一套经验：真正改变入口的技术，最后一定会触碰底层系统。
\end{importantbox}

\subsection{本章小结}

祝铭明的阿里前史训练了两种能力：判断平台入口何时切换，以及在大公司与创业公司之间选择适合自己的组织形态。这为后来的 AI/AR 双线埋下了基础。

